\documentclass[11pt,a4paper]{article}
\usepackage[utf8]{inputenc}
\usepackage[T1]{fontenc}
\usepackage{geometry}
\usepackage{graphicx}
\usepackage{booktabs}
\usepackage{enumitem}
\geometry{margin=1in}

\title{\textbf{Main Project 2: Simple Overview}\\
Code Organization, Models, and Data}
\author{}
\date{\today}

\begin{document}

\maketitle

\section{Project Overview}

This project analyzes macroeconomic regimes, forecasts economic variables, and builds asset allocation strategies. The code is organized into three main sections, each with a clear purpose.

\section{Data Used}

\subsection{Market Data}
\begin{itemize}
    \item \textbf{S\&P 500 Returns}: Monthly stock returns (1990-2025, 431 months)
    \item \textbf{3-Month Treasury Yield}: Converted to monthly bond returns
    \item \textbf{Equity Risk Premium (ERP)}: Stock return minus bond return
\end{itemize}

\subsection{Macro Variables}
Four standardized indices created from multiple economic indicators:
\begin{itemize}
    \item \textbf{Growth Factor}: Economic growth indicators
    \item \textbf{Inflation Factor}: Price level indicators  
    \item \textbf{Monetary Policy Factor}: Interest rates and policy measures
    \item \textbf{Market Volatility Factor}: Market volatility measures
\end{itemize}

\subsection{Sentiment Data}
Four sentiment scores extracted from news articles using AI:
\begin{itemize}
    \item Growth sentiment, Inflation sentiment, Policy sentiment, Volatility sentiment
\end{itemize}

\section{Code Organization}

The project follows a \textbf{modular structure} where each section has:
\begin{itemize}
    \item \texttt{preprocessing.py} - Loads and prepares data
    \item \texttt{model.py} or \texttt{*\_model.py} - Implements the main model
    \item \texttt{stats.py} - Statistical analysis and tests
    \item \texttt{plotting.py} - Creates visualizations
    \item \texttt{main.py} - Runs everything together
\end{itemize}

\section{Section 1: Macro Variable Analysis (s1\_macro\_vars)}

\subsection{s11\_full\_sample: Basic Regressions}
\textbf{What it does}: Tests how macro variables predict future stock returns (ERP).

\textbf{Model}: Simple linear regression
\[
\text{ERP}_{t+h} = \alpha + \beta_1 \text{Growth}_t + \beta_2 \text{Inflation}_t + \beta_3 \text{Policy}_t + \beta_4 \text{Volatility}_t + \varepsilon
\]

\textbf{Horizons}: 1, 3, 6, 12, 24 months ahead

\textbf{Outputs}: Which variables matter most, how strong the relationships are

\subsection{s12\_regimeness: Regime Detection}
\textbf{What it does}: Identifies different economic "regimes" or states.

\textbf{Three approaches}:
\begin{enumerate}
    \item \textbf{2×2 Simple Regimes}: Divides economy into 4 quadrants based on Growth (High/Low) × Inflation (High/Low)
    \begin{itemize}
        \item Goldilocks (High Growth, Low Inflation)
        \item Overheating (High Growth, High Inflation)
        \item Stagflation (Low Growth, High Inflation)
        \item Slowdown (Low Growth, Low Inflation)
    \end{itemize}
    
    \item \textbf{HMM (Hidden Markov Model)}: Uses statistical model to find regimes automatically
    \begin{itemize}
        \item Best model uses Growth + Policy (2 variables, 4 regimes)
        \item Provides probabilities for each regime at each time point
    \end{itemize}
    
    \item \textbf{Conditional Regressions}: Tests if macro-ERP relationships differ across regimes
\end{enumerate}

\subsection{s13\_extremeness: Extreme States Detection}
\textbf{What it does}: Identifies periods when macro variables are at extreme levels.

\textbf{Two methods}:
\begin{enumerate}
    \item \textbf{Isolation Forest}: Machine learning method to detect anomalies
    \item \textbf{PCA Distance}: Measures distance from "normal" in statistical space
\end{enumerate}

\textbf{Key finding}: Extreme macro states have 2× higher volatility in stock returns, making them useful for risk management.

\section{Section 2: Forecasting (s2\_forecasts)}

\textbf{Goal}: Forecast GDP growth and inflation 1, 3, and 6 months ahead.

\subsection{s21\_macro: TVP-VAR Model}
\textbf{What it is}: Time-Varying Parameter Vector Autoregression

\textbf{How it works}:
\begin{itemize}
    \item Uses all 4 macro variables together
    \item Allows relationships to change over time (time-varying)
    \item Implemented as rolling window VAR (practical approximation)
\end{itemize}

\textbf{Outputs}: Forecasts, time-varying coefficients, impulse response functions

\subsection{s22\_ml\_based: Machine Learning Models}

\subsubsection{XGBoost}
\textbf{What it is}: Gradient boosting machine learning model

\textbf{Two versions}:
\begin{itemize}
    \item \textbf{Macro-only}: Uses only macro variables
    \item \textbf{Macro + Sentiment}: Adds sentiment scores
\end{itemize}

\textbf{Features}: Includes ARIMA-like features (lags, differences, trends, volatility) to capture time series patterns

\subsubsection{LSTM}
\textbf{What it is}: Long Short-Term Memory neural network

\textbf{How it works}:
\begin{itemize}
    \item Uses sequences of past 12-24 months
    \item Learns patterns in the sequence
    \item Hybrid approach: Combines LSTM predictions with autoregressive terms
\end{itemize}

\textbf{Special features}: Custom loss function to prevent "flat" forecasts

\subsection{cross\_comparison: Model Comparison}
\textbf{What it does}: Compares all forecasting models

\textbf{Metrics}: RMSE, MAE for each model, variable, and horizon

\textbf{Outputs}: Performance tables, rankings, improvement charts

\section{Section 3: Asset Allocation Strategies (s3\_evaluation)}

\textbf{Goal}: Use regimes, extremeness, and forecasts to decide how much to invest in stocks vs bonds.

\subsection{Benchmark}
50/50 static allocation (50\% stocks, 50\% bonds), rebalanced monthly

\subsection{Strategy 1: Regime-Based}
\textbf{Idea}: Adjust stock allocation based on economic regime

\textbf{How}: 
\begin{itemize}
    \item Expansion regime → 70\% stocks
    \item Neutral regime → 50\% stocks  
    \item Risk-off regime → 30\% stocks
    \item Crisis regime → 20\% stocks
\end{itemize}

\textbf{Weight calculation}: Weighted average based on regime probabilities

\subsection{Strategy 2: Extremeness-Based}
\textbf{Idea}: Reduce stock allocation when macro conditions are extreme

\textbf{How}:
\begin{itemize}
    \item If extremeness > 90th percentile → 20\% stocks (risk-off)
    \item Otherwise → 60\% stocks (normal)
\end{itemize}

\textbf{Variant}: Combine with regime strategy for better performance

\subsection{Strategy 3: Forecast-Based}
\textbf{Idea}: Use forecasted stock returns (ERP) to adjust allocation

\textbf{Three methods}:
\begin{enumerate}
    \item \textbf{Threshold}: If forecast > 0 → 70\% stocks, else 30\%
    \item \textbf{Linear}: Map forecast to weight between 20\%-80\%
    \item \textbf{Regime-Conditioned}: Only act if forecast AND regime agree
\end{enumerate}

\textbf{Models tested}: TVP-VAR, XGBoost (macro), XGBoost (sentiment), LSTM

\subsection{Strategy 4: Ensemble}
\textbf{Idea}: Combine all signals (forecast + regime + extremeness)

\textbf{Rule}: 
\begin{itemize}
    \item Low weight if forecast < 0 OR extremeness high
    \item High weight if forecast > 0 AND regime is expansion
    \item Neutral otherwise
\end{itemize}

\subsection{Performance Metrics}
Each strategy evaluated on:
\begin{itemize}
    \item Annualized return
    \item Volatility (risk)
    \item Sharpe ratio (return per unit risk)
    \item Maximum drawdown (worst loss)
    \item Hit rate (prediction accuracy)
\end{itemize}

\section{Key Results Summary}

\subsection{Regime Detection}
\begin{itemize}
    \item HMM finds 4 meaningful economic regimes
    \item Best model uses Growth + Policy (2 variables)
    \item Regimes show different stock return patterns
\end{itemize}

\subsection{Extremeness}
\begin{itemize}
    \item Extreme macro states → 2× higher stock return volatility
    \item Useful for risk management (reducing exposure)
\end{itemize}

\subsection{Forecasting}
\begin{itemize}
    \item TVP-VAR performs well, especially for inflation
    \item Adding sentiment helps XGBoost for some horizons
    \item LSTM competitive with hybrid AR-LSTM approach
    \item ARIMA features improve both XGBoost and LSTM
\end{itemize}

\subsection{Strategies}
\begin{itemize}
    \item Regime-based: Smooth allocation adjustments
    \item Extremeness-based: Effective crash avoidance
    \item Forecast-based: Performance varies by model
    \item Ensemble: Typically best overall performance
\end{itemize}

\section{File Structure}

\begin{verbatim}
main_project2/
├── data/                    # All data files
│   ├── macro_final/        # Macro indices
│   ├── macro_processed/    # Market data (S&P500, bonds)
│   └── news_data/         # Sentiment scores
│
├── s1_macro_vars/         # Section 1: Macro analysis
│   ├── s11_full_sample/   # Basic regressions
│   ├── s12_regimeness/    # Regime detection
│   └── s13_extremeness/   # Extremeness analysis
│
├── s2_forecasts/          # Section 2: Forecasting
│   ├── s21_macro/         # TVP-VAR model
│   ├── s22_ml_based/      # XGBoost & LSTM
│   └── cross_comparison/  # Model comparison
│
└── s3_evaluation/         # Section 3: Strategies
    ├── strategy_regime.py
    ├── strategy_extremeness.py
    ├── strategy_forecast.py
    ├── strategy_ensemble.py
    └── main.py
\end{verbatim}

\section{How to Run}

Each section can be run independently:

\begin{itemize}
    \item \textbf{Section 1}: Run \texttt{main.py} in each subdirectory
    \item \textbf{Section 2}: Run \texttt{main.py} in s21\_macro, s22\_ml\_based, and cross\_comparison
    \item \textbf{Section 3}: Run \texttt{main.py} in s3\_evaluation
\end{itemize}

All results are saved in \texttt{results/} directories within each section.

\end{document}

